\documentclass[12pt,a4paper]{article}

% ---------------------------------------------------------
% PACKAGES
% ---------------------------------------------------------
\usepackage[a4paper,margin=1in]{geometry}
\usepackage{graphicx}
\usepackage{booktabs}
\usepackage{amsmath,amssymb}
\usepackage{float}
\usepackage{hyperref}
\usepackage{xcolor}
\usepackage{listings}
\usepackage{enumitem}
\usepackage{fancyhdr}
\usepackage{longtable}
\usepackage{array}
\usepackage{multirow}
\usepackage{caption}
\usepackage{subcaption}
\usepackage{siunitx}
\usepackage{microtype}

% ---------------------------------------------------------
% HYPERREF
% ---------------------------------------------------------
\hypersetup{
    colorlinks=true,
    linkcolor=blue,
    urlcolor=blue,
    citecolor=blue
}

% ---------------------------------------------------------
% PAGE STYLE
% ---------------------------------------------------------
\pagestyle{fancy}
\fancyhf{}
\lhead{ICS1512}
\rhead{Machine Learning Algorithms Laboratory}
\cfoot{\thepage}

% ---------------------------------------------------------
% CODE STYLE
% ---------------------------------------------------------
\lstset{
    language=Python,
    basicstyle=\ttfamily\small,
    numbers=left,
    numberstyle=\tiny,
    frame=single,
    breaklines=true,
    showstringspaces=false,
    columns=fullflexible,
    keepspaces=true
}

% ---------------------------------------------------------
% CUSTOM COMMAND FOR FIGURES
% ---------------------------------------------------------
% If a figure file exists, it is displayed.
% Otherwise, a placeholder is displayed so that the
% document still compiles.
\newcommand{\reportfigure}[4][0.75\linewidth]{
\begin{figure}[H]
    \centering
    \IfFileExists{#2}{
        \includegraphics[width=#1]{#2}
    }{
        \fbox{
            \parbox[c][5cm][c]{0.75\linewidth}{
                \centering
                \textbf{Figure Placeholder}\\[0.3cm]
                \texttt{#2}\\[0.2cm]
                Insert the corresponding figure generated
                from the notebook here.
            }
        }
    }
    \caption{#3}
    \label{fig:#4}
\end{figure}
}

% ---------------------------------------------------------
% DOCUMENT
% ---------------------------------------------------------
\begin{document}

\begin{tabular}{|p{4.2cm}|p{10.3cm}|}
\hline
\textbf{Experiment No.} & 3 \\ \hline

\textbf{Experiment Title} &
Regression Analysis using Linear and Regularized Models \\ \hline

\textbf{GitHub Repository} &
\url{https://github.com/RadeshL/Machine-Learning-Laboratory/blob/main/MLassign3.ipynb} \\ \hline

\textbf{Student Name} &
L Radesh \\ \hline

\textbf{Register Number} &
3122247001047 \\ \hline

\textbf{Faculty} &
Dr. Poreddy Ajay Kumar Reddy \\ \hline

\textbf{Academic Year} &
2026--2027 \\ \hline

\textbf{Semester} &
VI \\ \hline

\textbf{Submission Date} &
16 August, 2026 \\ \hline

\end{tabular}

\vspace{0.5cm}

% =========================================================
% 1. OBJECTIVE
% =========================================================

\section{Objective}

The objective of this experiment is to implement and compare linear and
regularized regression models for predicting a continuous target variable.
The experiment uses a real-world loan application dataset in which the
target variable is the loan amount sanctioned.

The specific objectives are:

\begin{enumerate}
    \item To perform exploratory data analysis on the loan application
    dataset.
    
    \item To identify and handle missing values and outliers.
    
    \item To preprocess numerical and categorical features using
    appropriate transformation, scaling, and encoding techniques.
    
    \item To implement Linear Regression as a baseline model.
    
    \item To implement Ridge, Lasso, and Elastic Net regression models.
    
    \item To tune regularization hyperparameters using GridSearchCV and
    RandomizedSearchCV with 5-fold cross-validation.
    
    \item To evaluate the regression models using Mean Absolute Error
    (MAE), Mean Squared Error (MSE), Root Mean Squared Error (RMSE),
    and $R^2$ score.
    
    \item To visualize predicted values, residuals, training and
    validation errors, and regression coefficients.
    
    \item To analyze the effect of regularization on model complexity,
    overfitting, underfitting, and the bias--variance trade-off.
\end{enumerate}

% =========================================================
% 2. PROBLEM STATEMENT
% =========================================================

\section{Problem Statement}

The problem is to predict the amount of loan that will be sanctioned to a
customer based on information associated with a loan application.

The dataset contains both numerical and categorical information related to
customers, their financial characteristics, employment information,
property information, credit history, loan requirements, and other
application attributes.

Let the input feature vector be represented by

\[
\mathbf{x} =
[x_1,x_2,\ldots,x_p]^T
\]

and let the target variable be

\[
y = \text{Loan Sanction Amount (USD)}.
\]

The objective is to learn a regression function

\[
\hat{y}=f(\mathbf{x})
\]

such that the predicted loan sanction amount is as close as possible to
the actual sanctioned amount.

Four regression models are considered:

\begin{enumerate}
    \item Linear Regression
    \item Ridge Regression
    \item Lasso Regression
    \item Elastic Net Regression
\end{enumerate}

The performance of the models is compared using regression error metrics
and the effect of regularization is analyzed.

% =========================================================
% 3. DATASET DESCRIPTION
% =========================================================

\section{Dataset Description}

The dataset used in this experiment is the \textit{Predict Loan Amount
Data} dataset referenced by the laboratory manual. It is a real-world
regression dataset containing numerical and categorical features related
to loan applications.

The target variable is \texttt{Loan Sanction Amount (USD)}.

\subsection{Dataset Summary}

\begin{longtable}{|p{6cm}|p{8cm}|}
\hline
\textbf{Property} & \textbf{Description} \\ \hline

Dataset Name &
Predict Loan Amount Data \\ \hline

Dataset Source &
Kaggle \\ \hline

Initial Number of Samples &
30,000 \\ \hline

Initial Number of Features/Columns &
24 \\ \hline

Target Variable &
Loan Sanction Amount (USD) \\ \hline

Final Number of Columns &
49 \\ \hline

Number of Duplicate Rows &
0 \\ \hline

Train-Test Split &
80\% training, 20\% testing \\ \hline

Random State &
42 \\ \hline

Cross-Validation &
5-Fold Cross-Validation \\ \hline

\end{longtable}

\subsection{Original Dataset Features}

The original dataset contains the following attributes:

\begin{itemize}
    \item Customer ID
    \item Name
    \item Gender
    \item Age
    \item Income (USD)
    \item Income Stability
    \item Profession
    \item Type of Employment
    \item Location
    \item Loan Amount Request (USD)
    \item Current Loan Expenses (USD)
    \item Expense Type 1
    \item Expense Type 2
    \item Dependents
    \item Credit Score
    \item No. of Defaults
    \item Has Active Credit Card
    \item Property ID
    \item Property Age
    \item Property Type
    \item Property Location
    \item Co-Applicant
    \item Property Price
    \item Loan Sanction Amount (USD)
\end{itemize}

\subsection{Target Variable}

The prediction target is:

\begin{center}
\textbf{Loan Sanction Amount (USD)}
\end{center}

This is a continuous numerical variable, making the problem a supervised
regression problem.

% =========================================================
% 4. EXPLORATORY DATA ANALYSIS
% =========================================================

\section{Exploratory Data Analysis}

Exploratory Data Analysis (EDA) was performed before model training to
understand the structure and quality of the dataset.

The initial dataset contains 30,000 rows and 24 columns. No duplicate rows
were detected.

\subsection{Dataset Overview}

The first five rows of the original dataset showed information about
customers, income, employment, loan requirements, credit score, property
details, and the loan amount sanctioned.

The target variable is directly available in the dataset as
\texttt{Loan Sanction Amount (USD)}.

\subsection{Missing Value Analysis}

The initial missing-value analysis identified missing observations in
several columns.

The main columns containing missing values were:

\begin{itemize}
    \item Gender
    \item Income (USD)
    \item Income Stability
    \item Type of Employment
    \item Current Loan Expenses (USD)
    \item Dependents
    \item Credit Score
    \item Has Active Credit Card
    \item Property Age
    \item Property Location
    \item Loan Sanction Amount (USD)
\end{itemize}

No column contained more than 50\% missing values. Therefore, no columns
were removed due to excessive missing data.

\reportfigure
    {MLassign3_3_1.png}
    {Missing-value distribution across the original dataset columns.}
    {missing-values}

\textbf{Inference:}

The missing-value analysis shows that missing observations are distributed
across several numerical and categorical attributes. No feature exceeded
the 50\% missing-value threshold, allowing all original features to be
retained at this stage. Numerical missing values were subsequently
replaced using their median values, while categorical missing values were
replaced using their mode. This prevents loss of a substantial number of
samples while maintaining a complete dataset for subsequent processing.

\subsection{Statistical and Feature Distribution Analysis}
These visualizations were used to inspect the distributions,
central tendencies, spread, skewness, and potential outliers in the
numerical features.

% =========================================================
% Numerical Feature Distribution Plots - Group 1
% =========================================================

\begin{figure}[H]
    \centering

    % -------------------------------------------------
    % Row 1
    % -------------------------------------------------
    \begin{subfigure}{0.44\linewidth}
        \centering
        \includegraphics[width=\linewidth]{MLassign3_3_3.png}
        \caption{Age}
    \end{subfigure}
    \hfill
    \begin{subfigure}{0.44\linewidth}
        \centering
        \includegraphics[width=\linewidth]{MLassign3_3_4.png}
        \caption{Income (USD)}
    \end{subfigure}

    \vspace{0.15cm}

    % -------------------------------------------------
    % Row 2
    % -------------------------------------------------
    \begin{subfigure}{0.44\linewidth}
        \centering
        \includegraphics[width=\linewidth]{MLassign3_3_5.png}
        \caption{Loan Amount Request (USD)}
    \end{subfigure}
    \hfill
    \begin{subfigure}{0.44\linewidth}
        \centering
        \includegraphics[width=\linewidth]{MLassign3_3_6.png}
        \caption{Current Loan Expenses (USD)}
    \end{subfigure}
    \label{fig:numerical-distributions-1}
\end{figure}


% =========================================================
% Numerical Feature Distribution Plots - Group 2
% =========================================================

\begin{figure}[H]
    \centering

    % -------------------------------------------------
    % Row 1
    % -------------------------------------------------
    \begin{subfigure}{0.44\linewidth}
        \centering
        \includegraphics[width=\linewidth]{MLassign3_3_7.png}
        \caption{Dependents}
    \end{subfigure}
    \hfill
    \begin{subfigure}{0.44\linewidth}
        \centering
        \includegraphics[width=\linewidth]{MLassign3_3_8.png}
        \caption{Credit Score}
    \end{subfigure}

    \vspace{0.15cm}

    % -------------------------------------------------
    % Row 2
    % -------------------------------------------------
    \begin{subfigure}{0.44\linewidth}
        \centering
        \includegraphics[width=\linewidth]{MLassign3_3_9.png}
        \caption{No. of Defaults}
    \end{subfigure}
    \hfill
    \begin{subfigure}{0.44\linewidth}
        \centering
        \includegraphics[width=\linewidth]{MLassign3_3_11.png}
        \caption{Property Age}
    \end{subfigure}
    \label{fig:numerical-distributions-2}
\end{figure}


% =========================================================
% Numerical Feature Distribution Plots - Group 3
% =========================================================

\begin{figure}[H]
    \centering

    % -------------------------------------------------
    % Row 1
    % -------------------------------------------------
    \begin{subfigure}{0.44\linewidth}
        \centering
        \includegraphics[width=\linewidth]{MLassign3_3_12.png}
        \caption{Property Type}
    \end{subfigure}
    \hfill
    \begin{subfigure}{0.44\linewidth}
        \centering
        \includegraphics[width=\linewidth]{MLassign3_3_13.png}
        \caption{Co-Applicant}
    \end{subfigure}

    \vspace{0.15cm}

    % -------------------------------------------------
    % Row 2
    % -------------------------------------------------
    \begin{subfigure}{0.44\linewidth}
        \centering
        \includegraphics[width=\linewidth]{MLassign3_3_14.png}
        \caption{Property Price}
    \end{subfigure}
    \hfill
    \begin{subfigure}{0.44\linewidth}
        \centering
        \includegraphics[width=\linewidth]{MLassign3_3_15.png}
        \caption{Loan Sanction Amount (USD)}
    \end{subfigure}
    \label{fig:numerical-distributions-3}
\end{figure}

\textbf{Inference:}

The histograms provide information about the distribution, central
tendency, and skewness of the numerical variables, while the corresponding
boxplots help identify potential outliers. The plots show that the
numerical features have different scales and distribution characteristics,
which makes feature standardization important before applying the
regression algorithms. The boxplots also reveal observations lying away
from the central range of some variables, motivating the use of the IQR
method for outlier detection and treatment. These preprocessing steps help
produce a more consistent feature representation for the subsequent
regression models.

\subsection{Categorical Feature Analysis}

Count plots were generated for categorical features with a manageable
number of categories. High-cardinality attributes such as Customer ID and
Name were excluded from categorical visualization and encoding because
they contain too many unique categories.

% =========================================================
% Categorical Feature Plots - Group 1
% =========================================================

\begin{figure}[H]
    \centering

    \begin{subfigure}{0.46\linewidth}
        \centering
        \includegraphics[width=\linewidth]{MLassign3_3_17.png}
        \caption{Gender}
    \end{subfigure}
    \hfill
    \begin{subfigure}{0.46\linewidth}
        \centering
        \includegraphics[width=\linewidth]{MLassign3_3_18.png}
        \caption{Income Stability}
    \end{subfigure}

    \vspace{0.25cm}

    \begin{subfigure}{0.46\linewidth}
        \centering
        \includegraphics[width=\linewidth]{MLassign3_3_19.png}
        \caption{Profession}
    \end{subfigure}
    \hfill
    \begin{subfigure}{0.46\linewidth}
        \centering
        \includegraphics[width=\linewidth]{MLassign3_3_20.png}
        \caption{Type of Employment}
    \end{subfigure}
    \label{fig:categorical-group1}
\end{figure}


% =========================================================
% Categorical Feature Plots - Group 2
% =========================================================

\begin{figure}[H]
    \centering

    \begin{subfigure}{0.46\linewidth}
        \centering
        \includegraphics[width=\linewidth]{MLassign3_3_21.png}
        \caption{Location}
    \end{subfigure}
    \hfill
    \begin{subfigure}{0.46\linewidth}
        \centering
        \includegraphics[width=\linewidth]{MLassign3_3_22.png}
        \caption{Expense Type 1}
    \end{subfigure}

    \vspace{0.25cm}

    \begin{subfigure}{0.46\linewidth}
        \centering
        \includegraphics[width=\linewidth]{MLassign3_3_23.png}
        \caption{Expense Type 2}
    \end{subfigure}
    \hfill
    \begin{subfigure}{0.46\linewidth}
        \centering
        \includegraphics[width=\linewidth]{MLassign3_3_24.png}
        \caption{Has Active Credit Card}
    \end{subfigure}
    \label{fig:categorical-group2}
\end{figure}


% =========================================================
% Categorical Feature Plots - Group 3
% =========================================================

\begin{figure}[H]
    \centering

    \begin{subfigure}{0.46\linewidth}
        \centering
        \includegraphics[width=\linewidth]{MLassign3_3_25.png}
        \caption{Property Location}
    \end{subfigure}
    \label{fig:categorical-group3}
\end{figure}


\textbf{Inference:}

The categorical plots show the distribution of observations across the
different categories. Some categorical variables contain only two
categories and are therefore suitable for binary encoding, while
variables with multiple categories require one-hot encoding. Customer ID
and Name were treated as high-cardinality attributes and were not
encoded.

% =========================================================
% 5. DATA PREPROCESSING
% =========================================================

\section{Data Preprocessing}

The preprocessing procedure consists of missing-value handling, outlier
treatment, numerical transformation, categorical encoding, and removal of
non-predictive identifier columns before model training.

\subsection{Missing Value Handling}

For numerical columns, missing values were replaced using the median:

\[
x_{\text{missing}} = \operatorname{median}(x).
\]

For categorical columns, missing values were replaced using the most
frequently occurring category:

\[
x_{\text{missing}} = \operatorname{mode}(x).
\]

This procedure resulted in zero remaining missing values after
preprocessing.

\subsection{Outlier Detection}

Outliers were detected using the Interquartile Range (IQR) method.

The first quartile and third quartile are represented by $Q_1$ and $Q_3$.
The IQR is

\[
IQR = Q_3-Q_1.
\]

The lower and upper limits were calculated as

\[
L = Q_1 - 1.5(IQR)
\]

and

\[
U = Q_3 + 1.5(IQR).
\]

Observations outside this interval were treated as outliers.

The implementation used a threshold of 15\% to distinguish between a
relatively small number of outliers and a more substantial proportion.

For variables with a small outlier percentage, outliers were replaced by
the median. For variables with a higher proportion of outliers, the values
were clipped to the IQR boundaries.

\subsection{Observed Outlier Treatment}

\begin{longtable}{|p{5.2cm}|p{2.5cm}|p{2.5cm}|p{4.5cm}|}
\hline
\textbf{Feature} & \textbf{Outliers} & \textbf{Percentage} &
\textbf{Treatment} \\ \hline

Age & 0 & 0\% & No treatment \\ \hline

Income (USD) & 1874 & 6.25\% &
Median replacement \\ \hline

Loan Amount Request (USD) & 752 & 2.51\% &
Median replacement \\ \hline

Current Loan Expenses (USD) & 874 & 2.91\% &
Median replacement \\ \hline

Dependents & 432 & 1.44\% &
Median replacement \\ \hline

Credit Score & 0 & 0\% &
No treatment \\ \hline

No. of Defaults & 5818 & 19.39\% &
IQR boundary clipping \\ \hline

Property ID & 0 & 0\% &
No treatment \\ \hline

Property Age & 1930 & 6.43\% &
Median replacement \\ \hline

Property Type & 0 & 0\% &
No treatment \\ \hline

Co-Applicant & 4484 & 14.95\% &
Median replacement \\ \hline

Property Price & 863 & 2.88\% &
Median replacement \\ \hline

Loan Sanction Amount (USD) & 453 & 1.51\% &
Median replacement \\ \hline

\end{longtable}

\subsection{Numerical Feature Scaling}

The skewness for each numerical feature was calculated.

The transformation rule was:

\begin{itemize}
    \item If $|\text{skewness}| > 1$ and the feature was non-negative,
    \texttt{log1p} transformation was considered.
    
    \item If $|\text{skewness}| > 1$ and negative values existed,
    StandardScaler was used.
    
    \item For features with $|\text{skewness}| \leq 1$, StandardScaler was
    applied.
\end{itemize}

In the actual dataset, the reported skewness values were within the
reasonable range for all listed numerical features. Therefore,
StandardScaler was applied to the numerical variables.

Standardization is given by

\[
z = \frac{x-\mu}{\sigma},
\]

where $\mu$ is the feature mean and $\sigma$ is its standard deviation.

\subsection{Categorical Encoding}

Binary categorical features were encoded using binary mappings.

For example, the notebook encoded:

\begin{itemize}
    \item Gender
    \item Income Stability
    \item Expense Type 1
    \item Expense Type 2
\end{itemize}

Categorical features with more than two categories were transformed using
one-hot encoding with the first category dropped.

The following variables were one-hot encoded:

\begin{itemize}
    \item Profession
    \item Type of Employment
    \item Location
    \item Has Active Credit Card
    \item Property Location
\end{itemize}

Customer ID and Name were not encoded because they have high cardinality.

\subsection{Final Dataset}

After preprocessing, the dataset contained:

\[
30,000 \text{ samples} \times 49 \text{ columns}.
\]

The identifier columns Customer ID and Name were subsequently excluded
from the model feature matrix.

The final model therefore operates on numerical and encoded categorical
predictors.

% =========================================================
% 6. TRAIN TEST SPLIT
% =========================================================

\section{Train-Test Split}

The dataset was divided into training and testing subsets using an
80:20 ratio.
Thus, approximately 24,000 observations were used for training and
6,000 observations were used for testing.

% =========================================================
% 7. MACHINE LEARNING ALGORITHMS
% =========================================================

\section{Machine Learning Algorithms}

Four regression algorithms were implemented:

\begin{enumerate}
    \item Linear Regression
    \item Ridge Regression
    \item Lasso Regression
    \item Elastic Net Regression
\end{enumerate}

\subsection{Linear Regression}

Linear Regression models the target variable as a linear combination of
the input features.

The mathematical formulation is

\[
\hat{y} = \beta_0 + \beta_1x_1+\beta_2x_2+\cdots+\beta_px_p.
\]

The model estimates the coefficients by minimizing the residual sum of
squares:

\[
\min_{\beta}
\sum_{i=1}^{n}
(y_i-\hat{y}_i)^2.
\]

Linear Regression provides a simple and interpretable baseline.

\subsubsection{Advantages}

\begin{itemize}
    \item Simple to implement.
    \item Computationally efficient.
    \item Easy to interpret.
    \item Provides a useful baseline for comparison.
\end{itemize}

\subsubsection{Limitations}

\begin{itemize}
    \item Sensitive to multicollinearity.
    \item Can overfit when many correlated features are present.
    \item Assumes a linear relationship between predictors and target.
\end{itemize}

\subsection{Ridge Regression}

Ridge Regression introduces an $L_2$ regularization penalty.

Its objective function is

\[
\min_{\beta}
\left[
\sum_{i=1}^{n}
(y_i-\hat{y}_i)^2
+
\alpha\sum_{j=1}^{p}\beta_j^2
\right].
\]

The regularization term penalizes large coefficient magnitudes.

A larger value of $\alpha$ produces stronger regularization.

\subsubsection{Advantages}

\begin{itemize}
    \item Reduces coefficient magnitudes.
    \item Helps reduce variance.
    \item Can improve generalization in the presence of correlated
    features.
\end{itemize}

\subsubsection{Limitations}

\begin{itemize}
    \item Does not generally force coefficients exactly to zero.
    \item Therefore, it does not perform explicit feature selection.
\end{itemize}

\subsection{Lasso Regression}

Lasso Regression uses an $L_1$ regularization penalty.

Its objective function is

\[
\min_{\beta}
\left[
\sum_{i=1}^{n}
(y_i-\hat{y}_i)^2
+
\alpha\sum_{j=1}^{p}|\beta_j|
\right].
\]

The $L_1$ penalty can drive some coefficients exactly to zero.

\subsubsection{Advantages}

\begin{itemize}
    \item Performs automatic feature selection.
    \item Produces sparse models.
    \item Can improve interpretability.
\end{itemize}

\subsubsection{Limitations}

\begin{itemize}
    \item Can discard useful correlated features.
    \item Performance is sensitive to the regularization parameter.
\end{itemize}

\subsection{Elastic Net Regression}

Elastic Net combines $L_1$ and $L_2$ regularization.

Its objective function can be represented as

\[
\min_{\beta}
\left[
\sum_{i=1}^{n}
(y_i-\hat{y}_i)^2
+
\alpha
\left(
\rho\sum_{j=1}^{p}|\beta_j|
+
(1-\rho)\sum_{j=1}^{p}\beta_j^2
\right)
\right],
\]

where $\rho$ is the L1 ratio.

When $\rho$ is close to 1, the model behaves more like Lasso, while lower
values increase the contribution of Ridge-style regularization.

\subsubsection{Advantages}

\begin{itemize}
    \item Combines feature selection and coefficient shrinkage.
    \item Can handle correlated predictors more effectively than pure
    Lasso.
    \item Provides a flexible balance between L1 and L2 penalties.
\end{itemize}

\subsubsection{Limitations}

\begin{itemize}
    \item Has more hyperparameters to tune.
    \item Can require greater computational time than basic Linear
    Regression.
\end{itemize}

% =========================================================
% 8. EXPERIMENTAL SETUP
% =========================================================

\section{Experimental Setup}

\begin{tabular}{ll}
\toprule
\textbf{Parameter} & \textbf{Configuration} \\
\midrule
Python Version & 3.13.3 \\
Notebook Kernel & Python 3 \\
Machine Learning Library & Scikit-learn \\
Data Processing & Pandas, NumPy \\
Visualization & Matplotlib, Seaborn \\
IDE & Jupyter Notebook \\
Operating System & Windows \\
Hardware & None \\
Random State & 42 \\
Train-Test Ratio & 80:20 \\
Cross-Validation & 5-Fold \\
\bottomrule
\end{tabular}

\vspace{0.4cm}

% =========================================================
% 9. HYPERPARAMETER SELECTION
% =========================================================

\section{Hyperparameter Selection}

The regularized regression models were tuned using both GridSearchCV and
RandomizedSearchCV.

The search was performed using 5-fold cross-validation and
\texttt{neg\_mean\_squared\_error} as the optimization score.

\subsection{Search Spaces}

\begin{longtable}{|p{4cm}|p{8cm}|}
\hline
\textbf{Model} & \textbf{Search Space} \\ \hline

Ridge &
$\alpha \in \{0.01,0.1,1,10,100\}$ \\ \hline

Lasso &
$\alpha \in \{0.001,0.01,0.1,1,10\}$ \\ \hline

Elastic Net &
$\alpha \in \{0.01,0.1,1,10\}$ and
$l1\_ratio \in \{0.2,0.5,0.8\}$ \\ \hline

\end{longtable}

\subsection{Grid Search Results}

\begin{longtable}{|p{3.5cm}|p{3.2cm}|p{3.2cm}|p{3cm}|}
\hline
\textbf{Model} & \textbf{Best Parameters} &
\textbf{Test RMSE} & \textbf{Test $R^2$} \\ \hline

Ridge &
$\alpha=100$ &
0.712176 &
0.490816 \\ \hline

Lasso &
$\alpha=0.001$ &
0.712226 &
0.490744 \\ \hline

Elastic Net &
$\alpha=0.01,\ l1\_ratio=0.2$ &
0.711781 &
0.491379 \\ \hline

\end{longtable}

The best-performing tuned configuration according to the test-set
$R^2$ was Elastic Net with

\[
\boxed{\alpha=0.01,\qquad l1\_ratio=0.2}.
\]

% =========================================================
% 10. RESULTS
% =========================================================

\section{Results}

\subsection{Baseline and Regularized Model Performance}

\begin{longtable}{|p{4cm}|c|c|c|c|c|}
\hline
\textbf{Model} &
\textbf{MAE} &
\textbf{MSE} &
\textbf{RMSE} &
\textbf{$R^2$} &
\textbf{Training Time (s)}
\\ \hline

Linear Regression &
0.500173 &
0.507597 &
0.712458 &
0.490412 &
0.072413
\\ \hline

Ridge Regression &
0.500169 &
0.507593 &
0.712455 &
0.490416 &
0.016880
\\ \hline

Lasso Regression &
0.523793 &
0.526109 &
0.725333 &
0.471827 &
0.074277
\\ \hline

Elastic Net &
0.513970 &
0.513546 &
0.716621 &
0.484439 &
0.091216
\\ \hline

\end{longtable}

\subsection{Tuned Model Performance}

\begin{longtable}{|p{4cm}|c|c|c|c|c|}
\hline
\textbf{Model} &
\textbf{MAE} &
\textbf{MSE} &
\textbf{RMSE} &
\textbf{$R^2$} &
\textbf{Training Time (s)}
\\ \hline

Grid Ridge &
0.500239 &
0.507194 &
0.712176 &
0.490816 &
0.538046
\\ \hline

Grid Lasso &
0.499948 &
0.507265 &
0.712226 &
0.490744 &
4.881348
\\ \hline

Grid Elastic Net &
0.500990 &
0.506633 &
0.711781 &
0.491379 &
5.851213
\\ \hline

\end{longtable}

\subsection{Randomized Search Performance}

\begin{longtable}{|p{4cm}|c|c|c|c|c|}
\hline
\textbf{Model} &
\textbf{Best Parameters} &
\textbf{MAE} &
\textbf{RMSE} &
\textbf{$R^2$} &
\textbf{Time (s)}
\\ \hline

Randomized Ridge &
$\alpha=100$ &
0.500239 &
0.712176 &
0.490816 &
0.512971
\\ \hline

Randomized Lasso &
$\alpha=0.001$ &
0.499948 &
0.712226 &
0.490744 &
4.906287
\\ \hline

Randomized Elastic Net &
$\alpha=0.01,\ l1\_ratio=0.2$ &
0.500990 &
0.711781 &
0.491379 &
5.545601
\\ \hline

\end{longtable}

An important observation is that GridSearchCV and RandomizedSearchCV
selected the same best parameters for all three regularized models.
Consequently, both approaches produced identical test-set performance in
this experiment.

% =========================================================
% 11. METRIC INTERPRETATION
% =========================================================

\section{Performance Metrics}

\subsection{Mean Absolute Error}

Mean Absolute Error is defined as

\[
MAE =
\frac{1}{n}
\sum_{i=1}^{n}
|y_i-\hat{y}_i|.
\]

It measures the average absolute difference between the actual and
predicted values.

Lower MAE indicates better prediction accuracy.

\subsection{Mean Squared Error}

Mean Squared Error is

\[
MSE =
\frac{1}{n}
\sum_{i=1}^{n}
(y_i-\hat{y}_i)^2.
\]

MSE penalizes larger errors more strongly than MAE.

\subsection{Root Mean Squared Error}

RMSE is

\[
RMSE = \sqrt{MSE}.
\]

It represents the square root of the average squared prediction error.

\subsection{$R^2$ Score}

The coefficient of determination is

\[
R^2 =
1-
\frac{
\sum_{i=1}^{n}(y_i-\hat{y}_i)^2
}{
\sum_{i=1}^{n}(y_i-\bar{y})^2
}.
\]

An $R^2$ value closer to 1 indicates that the model explains a greater
proportion of the variance in the target.

The best tuned model obtained an $R^2$ of approximately

\[
R^2=0.4914.
\]

This indicates that approximately 49.14\% of the variance in the
preprocessed target is explained by the tuned Elastic Net model on the
test set.

\subsection{Important Scale Consideration}

Applies StandardScaler to all numerical columns selected
during preprocessing, including the target variable
\texttt{Loan Sanction Amount (USD)}.

Therefore, the reported MAE, MSE, and RMSE values are measured on the
standardized target scale rather than directly in USD.

Consequently, the numerical error values should not be interpreted as
dollar amounts.

For a production-quality regression experiment, the preferred approach
would be to scale only the predictor variables and retain the original
target scale, or alternatively inverse-transform predictions before
reporting errors in USD.

% =========================================================
% 12. RESULT VISUALIZATION
% =========================================================

\section{Result Visualization}

\begin{enumerate}
    \item Predicted versus actual values
    \item Residual plots
    \item Training versus validation/test RMSE
    \item Coefficient comparison
\end{enumerate}

\subsection{Predicted versus Actual Values}

A predicted-versus-actual plot compares the model predictions with the
true target values.

The ideal relationship is represented by

\[
\hat{y}=y.
\]

Thus, points close to the diagonal line indicate accurate predictions.

\reportfigure
    {MLassign3_11_0.png}
    {Predicted versus actual values for Linear Regression.}
    {predicted-actual-linear}

\textbf{Inference:}

The predicted-versus-actual visualization shows the degree to which the
predicted loan sanction amounts follow the actual values. Points close to
the diagonal reference line indicate accurate predictions. The spread of
points around the line demonstrates that the linear model does not explain
all variation in the target. This is consistent with its test $R^2$ of
approximately 0.49.

\reportfigure
    {MLassign3_11_1.png}
    {Predicted versus actual values for tuned Ridge Regression.}
    {predicted-actual-ridge}

\textbf{Inference:}

The tuned Ridge predictions follow a pattern similar to the baseline
Linear Regression model. Regularization produces only a small numerical
improvement in the test metrics. The resulting $R^2$ is approximately
0.491, indicating a modest improvement in generalization.

\reportfigure
    {MLassign3_11_2.png}
    {Predicted versus actual values for tuned Lasso Regression.}
    {predicted-actual-lasso}

\textbf{Inference:}

The tuned Lasso model produces predictions that are broadly similar to the
other linear models. Its test MAE is slightly lower than that of tuned
Ridge and Elastic Net, although its RMSE and $R^2$ are marginally worse
than those of Elastic Net. The relatively small difference between the
models indicates that regularization does not dramatically change the
prediction structure for this dataset.

\reportfigure
    {MLassign3_11_3.png}
    {Predicted versus actual values for tuned Elastic Net Regression.}
    {predicted-actual-elastic}

\textbf{Inference:}

Elastic Net provides the strongest overall test-set performance among the
tuned models. It obtains the lowest RMSE and highest $R^2$ in the reported
results. The prediction scatter therefore provides visual support for the
numerical conclusion that Elastic Net is the preferred tuned model in this
experiment.

% ---------------------------------------------------------
% RESIDUAL PLOTS
% ---------------------------------------------------------

\subsection{Residual Analysis}

A residual is defined as

\[
e_i = y_i-\hat{y}_i.
\]

An ideal residual plot should show residuals distributed approximately
randomly around zero without a strong systematic pattern.

\reportfigure
    {MLassign3_11_4.png}
    {Residual plot for Linear Regression.}
    {residual-linear}

\textbf{Inference:}

The residual plot is used to determine whether systematic prediction errors
remain after fitting the regression model. Residuals distributed around
zero indicate that the model captures the main linear relationship. Any
visible structure or changing spread would indicate that some information
in the data remains unexplained by the model.

\reportfigure
    {MLassign3_11_5.png}
    {Residual plot for tuned Ridge Regression.}
    {residual-ridge}

\textbf{Inference:}

The tuned Ridge residual plot provides a diagnostic view of the effect of
$L_2$ regularization. A similar residual pattern to Linear Regression
indicates that Ridge has not substantially changed the underlying
prediction relationship. Its primary contribution is coefficient
shrinkage and modest variance control.

\reportfigure
    {MLassign3_11_6.png}
    {Residual plot for tuned Lasso Regression.}
    {residual-lasso}

\textbf{Inference:}

The Lasso residual plot shows the remaining prediction errors after
introducing $L_1$ regularization. Lasso encourages sparse coefficients,
which can simplify the model. However, its overall test $R^2$ remains
slightly below that of the tuned Elastic Net model.

\reportfigure
    {MLassign3_11_7.png}
    {Residual plot for tuned Elastic Net Regression.}
    {residual-elastic}

\textbf{Inference:}

The Elastic Net residual plot is used to assess whether the combined
regularization approach leaves systematic patterns in the errors. The
model's relatively low RMSE and highest test $R^2$ among the tuned models
suggest that its balance between coefficient shrinkage and sparsity is
effective for this dataset.

% ---------------------------------------------------------
% TRAINING VS VALIDATION
% ---------------------------------------------------------

\subsection{Training Error versus Validation/Test Error}

Compared training RMSE and test-set RMSE for the four
models.

\reportfigure
    {MLassign3_11_8.png}
    {Comparison of training RMSE and validation/test RMSE across models.}
    {training-validation}

\textbf{Inference:}

The training-error versus validation/test-error plot is useful for
detecting possible overfitting and underfitting. A very large difference
between training and validation error would indicate that the model is
fitting the training data substantially better than unseen data. The
regularized models provide a mechanism for controlling this gap by
penalizing excessive coefficient magnitudes. In this implementation, the
plot compares the training error with the held-out test error rather than
a separately stored cross-validation validation-error curve.

% ---------------------------------------------------------
% COEFFICIENT COMPARISON
% ---------------------------------------------------------

\subsection{Coefficient Comparison}

The regression coefficient for a feature represents its contribution to
the predicted target while holding the other features constant.

\reportfigure
    {MLassign3_11_9.png}
    {Coefficient comparison among Linear, Ridge, Lasso, and Elastic Net models.}
    {coefficient-comparison}

\textbf{Inference:}

The coefficient comparison illustrates the effect of regularization on
model parameters. Ridge reduces coefficient magnitudes through its
$L_2$ penalty, while Lasso can drive coefficients exactly to zero because
of its $L_1$ penalty. Elastic Net combines both effects. Therefore, the
coefficient plot provides a direct visualization of the trade-off between
model complexity and feature contribution.

% =========================================================
% 13. OVERFITTING AND UNDERFITTING
% =========================================================

\section{Overfitting and Underfitting Analysis}

Overfitting occurs when a model learns the training data too closely and
performs significantly worse on unseen data. Underfitting occurs when the
model is too simple to capture the underlying relationships.

\subsection{Linear Regression}

Linear Regression provides the baseline performance.

Its test-set performance was:

\[
MAE=0.500173
\]

\[
RMSE=0.712458
\]

\[
R^2=0.490412.
\]

The baseline model provides a useful reference for evaluating whether
regularization produces improvements.

\subsection{Effect of Ridge Regularization}

The baseline Ridge model has nearly identical performance to Linear
Regression.

The tuned Ridge model improves the test $R^2$ from approximately

\[
0.490412
\]

to

\[
0.490816.
\]

The improvement is small, indicating that Ridge regularization does not
dramatically change the predictive capability of the model.

\subsection{Effect of Lasso Regularization}

The initial Lasso model performs worse than Linear Regression:

\[
R^2=0.471827.
\]

After hyperparameter tuning, the Lasso model improves substantially:

\[
R^2=0.490744.
\]

This demonstrates the importance of selecting an appropriate
regularization strength.

\subsection{Effect of Elastic Net Regularization}

The initial Elastic Net configuration achieves

\[
R^2=0.484439.
\]

After tuning, its performance improves to

\[
R^2=0.491379.
\]

Thus, hyperparameter optimization produces the strongest test-set
performance among the evaluated models.

% =========================================================
% 14. BIAS VARIANCE
% =========================================================

\section{Bias--Variance Analysis}

The bias--variance trade-off describes the relationship between model
complexity, systematic error, and sensitivity to the training data.

\subsection{Linear Regression}

Linear Regression has relatively low model complexity and therefore
provides a useful baseline with relatively controlled variance. However,
its linear structure can introduce bias when the underlying relationship
between the features and target is more complex.

\subsection{Ridge Regression}

Ridge adds an $L_2$ penalty that shrinks coefficients toward zero.

Increasing $\alpha$ increases the strength of the penalty and generally
reduces model variance at the cost of potentially increasing bias.

The tuned Ridge model selected

\[
\alpha=100.
\]

This relatively large regularization value indicates that strong
coefficient shrinkage produced the best result within the tested search
space.

\subsection{Lasso Regression}

Lasso applies an $L_1$ penalty and can produce sparse coefficient vectors.

The tuned Lasso model selected

\[
\alpha=0.001.
\]

The small value suggests that strong L1 regularization was not preferred
within the tested search space.

The model can nevertheless perform feature selection by shrinking some
coefficients to zero.

\subsection{Elastic Net}

Elastic Net provides a compromise between L1 and L2 regularization.

The selected configuration was

\[
\alpha=0.01
\]

and

\[
l1\_ratio=0.2.
\]

A value of $0.2$ indicates a stronger contribution from the L2 component
than the L1 component. This configuration achieved the highest test
$R^2$ among the tuned models.

% =========================================================
% 15. GRID SEARCH VS RANDOM SEARCH
% =========================================================

\section{Grid Search and Randomized Search Comparison}

Both GridSearchCV and RandomizedSearchCV were implemented for Ridge,
Lasso, and Elastic Net.

The important observation is that both methods selected identical best
parameters:

\begin{longtable}{|p{3.5cm}|p{4cm}|p{4cm}|}
\hline
\textbf{Model} & \textbf{Grid Search} & \textbf{Randomized Search}
\\ \hline

Ridge &
$\alpha=100$ &
$\alpha=100$ \\ \hline

Lasso &
$\alpha=0.001$ &
$\alpha=0.001$ \\ \hline

Elastic Net &
$\alpha=0.01,\ l1\_ratio=0.2$ &
$\alpha=0.01,\ l1\_ratio=0.2$ \\ \hline

\end{longtable}

Since the randomized searches sampled the complete available search
spaces in this particular experiment, they produced the same best
configurations as GridSearchCV.

The measured execution times were approximately:

\begin{itemize}
    \item Ridge Grid Search: 0.538 seconds
    \item Ridge Randomized Search: 0.513 seconds
    \item Lasso Grid Search: 4.881 seconds
    \item Lasso Randomized Search: 4.906 seconds
    \item Elastic Net Grid Search: 5.851 seconds
    \item Elastic Net Randomized Search: 5.546 seconds
\end{itemize}

The small timing differences demonstrate that randomized search can have
a computational advantage when fewer configurations are sampled, although
the difference depends on the search space and number of iterations.

% =========================================================
% 16. DISCUSSION
% =========================================================

\section{Discussion}

The experiment demonstrates the implementation and comparison of four
linear regression approaches for predicting loan sanction amounts.

The baseline Linear Regression model achieved an $R^2$ score of
approximately 0.4904. Ridge produced an almost identical baseline result,
while the initial Lasso and Elastic Net configurations performed slightly
worse.

Hyperparameter tuning improved the performance of all regularized models.
The best Ridge parameter was $\alpha=100$, while Lasso selected
$\alpha=0.001$. Elastic Net selected $\alpha=0.01$ and
$l1\_ratio=0.2$.

Among the tuned models, Elastic Net achieved the best test performance:

\[
MAE=0.500990,
\]

\[
MSE=0.506633,
\]

\[
RMSE=0.711781,
\]

and

\[
R^2=0.491379.
\]

The improvement over baseline Linear Regression is relatively small.
Therefore, regularization improves the result, but it does not completely
solve the limitations of a linear model for this dataset.

\subsection{Strengths}

The experiment has several strengths:

\begin{itemize}
    \item A complete preprocessing procedure was implemented.
    \item Both numerical and categorical variables were handled.
    \item Missing values were systematically imputed.
    \item Outliers were explicitly detected and treated.
    \item Four regression algorithms were compared.
    \item Both GridSearchCV and RandomizedSearchCV were implemented.
    \item Five-fold cross-validation was used for hyperparameter selection.
    \item Multiple regression metrics were reported.
    \item Model predictions, residuals, errors, and coefficients were
    visualized.
\end{itemize}

\subsection{Limitations and Improvements}

All four models considered are linear regression models. If the underlying
relationship between loan characteristics and sanctioned amount is
nonlinear, models such as Random Forest Regression, Gradient Boosting,
XGBoost, or other nonlinear approaches could potentially provide better
performance.

% =========================================================
% 17. OBSERVATIONS
% =========================================================

\section{Observations}

The major observations from the experiment are:

\begin{enumerate}
    \item The dataset contains 30,000 loan application records and
    multiple numerical and categorical attributes.
    
    \item Missing values were present in several columns, but no column
    exceeded the 50\% missing-value threshold.
    
    \item Numerical missing values were handled using median imputation,
    while categorical missing values were handled using mode imputation.
    
    \item Outliers were identified using the IQR method.
    
    \item Numerical features were standardized using StandardScaler.
    
    \item Binary categorical variables were binary encoded, while
    multi-category variables were one-hot encoded.
    
    \item The final preprocessed dataset contained 49 columns.
    
    \item Linear Regression achieved an $R^2$ score of approximately
    0.4904.
    
    \item The initial Lasso and Elastic Net models performed worse than
    the baseline Linear Regression model.
    
    \item Hyperparameter tuning improved the performance of the
    regularized models.
    
    \item GridSearchCV and RandomizedSearchCV selected identical best
    hyperparameters for all three regularized models.
    
    \item Elastic Net achieved the highest test $R^2$ of approximately
    0.4914.
    
    \item Elastic Net also achieved the lowest test RMSE among the tuned
    models.
    
    \item The improvement from regularization was relatively small,
    suggesting that the predictive capability of the tested linear
    models is limited by the feature-target relationship.
\end{enumerate}

% =========================================================
% 18. CONCLUSION
% =========================================================

\section{Conclusion}

This experiment implemented Linear Regression, Ridge Regression, Lasso
Regression, and Elastic Net Regression for predicting loan sanction
amounts using a real-world loan application dataset.

The dataset was first explored and preprocessed through missing-value
imputation, IQR-based outlier handling, numerical standardization, and
categorical encoding. The resulting dataset contained 30,000 samples and
49 processed columns, from which the identifier columns Customer ID and
Name were excluded before model training.

Linear Regression was used as the baseline model. Ridge, Lasso, and
Elastic Net were subsequently trained to investigate the effect of
regularization.

Hyperparameter tuning was performed using both GridSearchCV and
RandomizedSearchCV with five-fold cross-validation. The best parameters
were $\alpha=100$ for Ridge, $\alpha=0.001$ for Lasso, and
$\alpha=0.01$, $l1\_ratio=0.2$ for Elastic Net.

The tuned Elastic Net model achieved the best test-set performance with

\[
MAE=0.500990,
\]

\[
MSE=0.506633,
\]

\[
RMSE=0.711781,
\]

and

\[
R^2=0.491379.
\]

Thus, Elastic Net is selected as the best model among the evaluated
configurations.

However, the improvement over the baseline Linear Regression model is
small. This indicates that regularization provides only a modest benefit
for this particular feature representation. More advanced nonlinear
regression methods, improved feature engineering, and a leakage-free
preprocessing pipeline could be explored to improve predictive
performance.

Finally, the experiment demonstrates the importance of model
regularization, hyperparameter tuning, cross-validation, and error
analysis when developing regression models for real-world prediction
tasks.

% =========================================================
% 19. REFERENCES
% =========================================================

\section{References}

\begin{enumerate}[label={[}\arabic*{]}]

    \item F. Pedregosa et al.,
    ``Scikit-learn: Machine Learning in Python,''
    \textit{Journal of Machine Learning Research},
    vol. 12, pp. 2825--2830, 2011.

    \item Scikit-learn Documentation,
    ``Linear Models.''

    \item Scikit-learn Documentation,
    ``Tuning the Hyper-Parameters of an Estimator.''

    \item Kaggle,
    ``Predict Loan Amount Data,'' dataset used for regression analysis.

\end{enumerate}

\end{document}